% Generated draft: requires mathematical and source review.
\documentclass[11pt]{article}
\usepackage[margin=25mm]{geometry}
\usepackage{fontspec}
\setmainfont{Latin Modern Roman}
\usepackage{amsmath,amssymb,mathtools,unicode-math}
\setmathfont{Latin Modern Math}
\usepackage{enumitem,needspace,xurl,hyperref}
\setlength{\parindent}{0pt}
\setlength{\parskip}{5pt}
\setlength{\emergencystretch}{3em}
\newcommand{\N}{\mathbb{N}}
\newcommand{\Z}{\mathbb{Z}}
\newcommand{\Q}{\mathbb{Q}}
\newcommand{\R}{\mathbb{R}}
\newcommand{\C}{\mathbb{C}}
\newcommand{\F}{\mathbb{F}}
\newcommand{\A}{\mathbb{A}}
\newcommand{\PP}{\mathbb{P}}
\newcommand{\E}{\mathbb{E}}
\newcommand{\eps}{\varepsilon}
\newcommand{\dd}{\,\mathrm{d}}
\begin{document}
\section*{UnsolvedMath ID: 700526}
\textbf{Source problem code:} catalogue entry
\subsection{Definitions and mathematical statement}
A group is called countable if its underlying set is finite or countably infinite. A group is periodic if every element has finite order. A group is locally finite, locally soluble, or locally nilpotent if every finitely generated subgroup is, respectively, finite, soluble, or nilpotent. Here a group is soluble if its derived series reaches the trivial subgroup after finitely many steps, and nilpotent if its lower central series reaches the trivial subgroup after finitely many steps. The required number of steps may depend on the finitely generated subgroup.

For a class \(\mathcal C\) of countable groups, call \(U\) universal for \(\mathcal C\) under embeddings if
\[
U\in\mathcal C
\quad\text{and}\quad
\forall H\in\mathcal C\;\exists\,\iota:H\hookrightarrow U,
\]
where \(\iota\) is an injective group homomorphism. Thus the source's word ``largest'' means containing an isomorphic copy of every member of the class; it imposes no uniqueness requirement.

The supplied source states as background that the class of countable locally finite groups has such a universal member. B. Hartley's problem asks which other classes of countable groups have a universal member under embeddings. In particular, determine separately whether each of the following assertions holds:
\begin{enumerate}
\item There exists a countable periodic locally soluble group \(U_{\mathrm s}\) such that every countable periodic locally soluble group embeds into \(U_{\mathrm s}\).
\item There exists a countable periodic locally nilpotent group \(U_{\mathrm n}\) such that every countable periodic locally nilpotent group embeds into \(U_{\mathrm n}\).
\end{enumerate}
The general request does not specify a collection of admissible classes or a required form of classification; the two displayed existence questions are its explicitly identified targets.
\subsection{Short English statement}
B. Hartley asks which classes of countable groups contain a countable group into which every group in the same class embeds. The two specific questions concern periodic locally soluble groups and periodic locally nilpotent groups. The supplied record labels the problem open; this draft does not independently verify that status.
\subsection{Sources}
\begin{itemize}
\item Local UnsolvedMath registry: supplied complete source record, permanent ID \(700526\), exact source problem code \texttt{catalogue entry}, title \emph{Kourovka Notebook Problem 8.78}; upstream revision \texttt{d371ee917ac9647fbb06cd6fc529deec6f3dda82}.
\item \emph{The Kourovka Notebook (Unsolved Problems in Group Theory)}, problem 8.78, attributed to B. Hartley, 8th issue (1982). The supplied background identifies the 21st issue edition (2026), \texttt{arXiv:1401.0300v48}. The supplied collection description names E. I. Khukhro and V. D. Mazurov as editors.
\item The existence assertion for countable locally finite groups and the label ``open'' are supplied source information. The cited edition and the mathematical status have not been independently checked for this draft.
\end{itemize}
\end{document}
